% from main.tex
%%%%%%%%%%  scan idx 544  |  folio 531  %%%%%%%%%%
\section*{2. Polynômes de Snapper}

\subsection*{2.1}
Rappelons d'abord la définition et quelques propriétés élémentaires des applications polynômes :

\textbf{Définition 2.1.1.} Soit $f$ une application d'un groupe abélien $A$ dans un groupe abélien $B$. On dit que $f$ est une application polynôme, s'il existe, pour toute suite $x_1,\ldots,x_r$ ($r\geq1$) d'éléments de $A$, une famille $(y(x_1,\ldots,x_r;k_1,\ldots,k_r))_{k_1,\ldots,k_r\geq0}$ d'éléments de $B$, presque tous nuls, et satisfaisant à la condition suivante : Pour toute suite $n_1,\ldots,n_r$ d'entiers, on a
\[
f(n_1x_1+\cdots+n_rx_r)=
\sum_{k_1,\ldots,k_r\geq0}(n_1^{k_1}\cdots n_r^{k_r})
y(x_1,\ldots,x_r;k_1,\ldots,k_r).
\]
On dit que $f$ est de degré $\leq m$, si quels que soient $r\in\mathbb Z$ et $x_1,\ldots,x_r\in A$, on a
\[
y(x_1,\ldots,x_r;k_1,\ldots,k_r)=0
\]
pour $k_1+\cdots+k_r>m$.

\textbf{Lemme 2.1.2.} Tout homomorphisme de groupes abéliens est une application polynôme de degré $\leq1$. Si $h=f\circ g$, où $f$ est une application polynôme de degré $p$ et $g$ une application polynôme de degré $q$, alors $h$ est une application polynôme de degré $\leq pq$.

Démonstration triviale.

\textbf{Lemme 2.1.3.} Soit
\[
0\longrightarrow G'\longrightarrow G\xrightarrow{j}G''\longrightarrow0
\]
une suite exacte de groupes abéliens telle que $G'$ est divisible par un
